\documentclass[11pt]{article}

\usepackage{amsmath}
\usepackage{amssymb}
\usepackage[margin=1in]{geometry}

\title{Effective Core Potential Integrals in PyFock}
\author{PyFock implementation notes}
\date{}

\begin{document}
\maketitle

\section{Semilocal ECP Operator}

The scalar semilocal effective core potential on center \(C\) is written as
\begin{equation}
  \hat U_C =
  U_L(r_C) +
  \sum_{\ell=0}^{L-1}\sum_{m=-\ell}^{\ell}
  |\ell m\rangle
  \left[ U_\ell(r_C)-U_L(r_C) \right]
  \langle \ell m | ,
\end{equation}
where \(r_C = |\mathbf r-\mathbf C|\), \(L\) is the local channel, and
\(|\ell m\rangle\) denotes a normalized real spherical harmonic projector.
The TURBOMOLE/def2 data in PyFock are parsed as Gaussian radial expansions
\begin{equation}
  U_\ell(r) = \sum_k c_{\ell k} r^{n_{\ell k}-2}
  \exp(-\alpha_{\ell k} r^2).
\end{equation}
For the def2 Cd ECP all tabulated radial powers are \(n=2\), so the radial
terms used by the local analytic path are ordinary Gaussians
\(c_k\exp(-\alpha_k r^2)\).

\section{Local Channel}

For Cartesian Gaussian primitives
\begin{align}
  \chi_a(\mathbf r) &=
  N_a (x-A_x)^{l_a}(y-A_y)^{m_a}(z-A_z)^{n_a}
  \exp[-\alpha|\mathbf r-\mathbf A|^2],\\
  \chi_b(\mathbf r) &=
  N_b (x-B_x)^{l_b}(y-B_y)^{m_b}(z-B_z)^{n_b}
  \exp[-\beta|\mathbf r-\mathbf B|^2],
\end{align}
and a local ECP Gaussian \(c\exp[-\gamma|\mathbf r-\mathbf C|^2]\), the integral is
\begin{equation}
  V^{\mathrm{loc}}_{ab,k}
  =
  c N_a N_b
  \exp\left[-\alpha A^2-\beta B^2-\gamma C^2
  + p P^2\right]
  S_x S_y S_z,
\end{equation}
with
\begin{equation}
  p=\alpha+\beta+\gamma,\qquad
  \mathbf P = \frac{\alpha\mathbf A+\beta\mathbf B+\gamma\mathbf C}{p}.
\end{equation}
The one-dimensional Hermite overlap factors are
\begin{equation}
  S_x =
  \int_{-\infty}^{\infty}
  (x-A_x)^{l_a}(x-B_x)^{l_b}
  \exp[-p(x-P_x)^2]\,dx ,
\end{equation}
and similarly for \(S_y\) and \(S_z\). PyFock evaluates these with the same
\texttt{calcS} helper used by the overlap integrals, then sums over primitive
pairs and ECP primitives.

\section{Projector Channels}

For a projector channel \(\ell\), the matrix contribution is
\begin{equation}
  V^{\ell}_{ab}
  =
  \sum_m \int_0^\infty r^2
  \left[U_\ell(r)-U_L(r)\right]
  q_{a\ell m}(r)\,q_{b\ell m}(r)\,dr ,
\end{equation}
where
\begin{equation}
  q_{a\ell m}(r)
  =
  \int_{S^2}
  \chi_a(\mathbf C+r\Omega)\,Y_{\ell m}(\Omega)\,d\Omega .
\end{equation}
The implementation in \texttt{ecp\_mat\_symm.py} uses normalized real spherical harmonics and forms these
projector amplitudes for all AOs at a radial node as a dense matrix product.
The angular integral is evaluated on a product grid that is exact for the low
angular momenta appearing in the def2 Cd basis at practical grid sizes. The
radial integral uses the same Gauss--Chebyshev mapping used in established ECP
implementations, but implemented natively in Python/Numba and without PySCF,
libcint, or libint calls.

\section{Analytical Projector Series}

The implementation in \texttt{ecp\_mat\_symm\_analytical.py} removes the
spatial quadrature from the projector channels. By the addition theorem,
\begin{equation}
  \sum_m Y_{\ell m}(\Omega)Y_{\ell m}(\Omega')
  =
  \frac{2\ell+1}{4\pi}P_\ell(\Omega\cdot\Omega'),
\end{equation}
so the two angular projections can be written as moments of monomials in
\(\Omega\) and \(\Omega'\). For an AO centered at \(\mathbf A\), define
\(\mathbf D=\mathbf A-\mathbf C\). Around the ECP center,
\begin{equation}
  \chi_a(\mathbf C+r\Omega)
  =
  N_a (\mathbf r_\Omega-\mathbf D)^{\mathbf l_a}
  \exp[-\alpha(r^2+D^2)]
  \exp[2\alpha r\,\mathbf D\cdot\Omega],
\end{equation}
where \(\mathbf r_\Omega=r\Omega\). The off-center exponential is expanded as
\begin{equation}
  \exp[2\alpha r\,\mathbf D\cdot\Omega]
  =
  \sum_{t_x,t_y,t_z\ge 0}
  \frac{(2\alpha r)^{t_x+t_y+t_z}
  D_x^{t_x}D_y^{t_y}D_z^{t_z}
  \Omega_x^{t_x}\Omega_y^{t_y}\Omega_z^{t_z}}
  {t_x!t_y!t_z!}.
\end{equation}
Every retained term has an exact angular integral over the unit sphere,
\begin{equation}
  \int_{S^2}\Omega_x^{2a}\Omega_y^{2b}\Omega_z^{2c}\,d\Omega
  =
  4\pi
  \frac{(2a-1)!!(2b-1)!!(2c-1)!!}
       {(2a+2b+2c+1)!!},
\end{equation}
with odd powers vanishing. The remaining radial integral is also closed form:
\begin{equation}
  \int_0^\infty r^q \exp[-p r^2]\,dr
  =
  \frac{1}{2}\Gamma\left(\frac{q+1}{2}\right)
  p^{-(q+1)/2}.
\end{equation}
Thus the analytical code performs no radial or angular quadrature. For
off-center AO/ECP pairs the exact expression is an infinite analytical series;
the \texttt{series\_order} parameter controls the truncation. Centered
AO/ECP terms are exact at order zero. For Cd\(_2\)/def2-TZVP at 2.98
\AA, order 12 reaches about \(10^{-8}\) maximum matrix-element error versus
PySCF, and order 20 reaches about \(10^{-11}\).

\section{Molecular Bookkeeping}

When an ECP block is found, PyFock stores the atom index, local channel, core
electron count, and projector channels on the \texttt{Basis} object. The
molecule's active nuclear charge is reduced by \(n_{\mathrm{core}}\), and the
active electron count is reduced by the same amount. Thus a neutral Cd atom with
a 28-electron def2 ECP has active charge 20 and contributes 20 valence
electrons. The ordinary nuclear attraction and nuclear repulsion terms then use
the valence charges, while \(\hat U_C\) supplies the core replacement operator.

\section{Validation}

For Cd/def2-TZVP, PyFock's Cartesian functions are internally normalized before
the Cartesian-to-spherical transformation. Therefore direct Cartesian matrices
do not numerically match PySCF's unnormalized Cartesian integral convention.
After transforming to the spherical AO basis used in the benchmark, the native
PyFock ECP matrix agrees with PySCF's \texttt{ECPscalar\_sph} matrix to
roundoff for an isolated Cd atom.

\begin{thebibliography}{9}
\bibitem{mclean1976}
A. D. McLean and M. Yoshimine, ``Theory of molecular integrals over Gaussian
functions,'' \emph{J. Chem. Phys.} \textbf{65}, 3826 (1976).

\bibitem{mitin1999}
A. V. Mitin, C. van Wüllen, and G. Hirsch, ``Two-component relativistic
effective core potentials,'' \emph{J. Chem. Phys.} \textbf{111}, 8778 (1999).

\bibitem{mura1996}
M. E. Mura and P. J. Knowles, ``Improved radial grids for quadrature in
molecular density-functional calculations,'' \emph{J. Chem. Phys.}
\textbf{104}, 9848 (1996).
\end{thebibliography}

\end{document}
